\documentclass[11pt]{article}
\usepackage[a4paper,margin=26mm]{geometry}
\usepackage{amsmath,amssymb}
\usepackage[hidelinks]{hyperref}
\setlength{\emergencystretch}{3em}
\title{Conjecture 00000004271: the rational divisible group obstructs approximation}
\author{lizaixi01, with OpenAI Codex assistance}
\date{23 October 2026}
\begin{document}
\maketitle
\section{The approximation obstruction}
Take the actual additive group $A=(\mathbf Q,+)$. It is divisible:
for every positive integer $n$ and $x\in\mathbf Q$, the element $x/n$
satisfies $n(x/n)=x$. The map $D(x)=2x$ is a nonidentity additive
automorphism, with inverse $x\mapsto x/2$.

On this group the identity is the only automorphism satisfying any
of the following three finite-support conditions:
\begin{enumerate}
\item only finitely many elements are moved;
\item the image $(f-\mathrm{id})(A)$ is finite;
\item $f$ acts identically on a subgroup of finite index.
\end{enumerate}
The first is the literal finite support of a permutation. The third
is the standard group-theoretic finitary condition; its equivalence
with the second is stated in Dardano and Rinauro,
\href{https://arxiv.org/abs/1403.4193v4}
{\emph{The group of inertial automorphisms of an abelian group}},
arXiv:1403.4193v4, Introduction, pp.1--2. We use that exact preprint
title and version. Lean proves all these claims for $\mathbf Q$,
including the finite-index bridge, instead of importing the quoted
equivalence as an assumption.

Thus no choice of a genuine separating metric can allow these
automorphisms to approximate $D$. The source does not specify a
formula for its ``subgroup-index metric.'' We avoid choosing a new
formula: the formal obstruction holds for \emph{every} true metric,
both for uniform approximation on $A$ and for approximation in its
automorphism group. This disproves the approximation necessary clause,
so the additional proposed optimal exponent cannot rescue the stated
conjunction.

\section{Actual finite-support and finite-index calculations}
Every additive rational automorphism satisfies
\[
 f(x)=f(1)x.
\]
Indeed additive homomorphisms between rational vector spaces commute
with rational scalar multiplication. If $f(1)\ne1$, then
\[
 (f-\mathrm{id})(x)=(f(1)-1)x
\]
is injective on the infinite set $\mathbf Q$; its image is therefore
infinite. Hence condition 2 forces $f(1)=1$ and $f=\mathrm{id}$.
Condition 1 implies condition 2 because the defect image is contained
in zero together with the image of the finite moved-element set.

\newpage
For condition 3 define the genuine additive subgroup
\[
 F=\ker(f-\mathrm{id})=\{x\in\mathbf Q:f(x)=x\}.
\]
If $f$ fixes some finite-index subgroup $H$ pointwise, then $H\le F$,
so $F$ has finite index as well. Conversely $F$ itself is such an $H$.
This equivalence is proved using Mathlib's actual
\texttt{AddSubgroup.FiniteIndex} and subgroup inclusion.

Let $N=[\mathbf Q:F]$, a positive natural number when $F$ has finite
index. The standard quotient-group calculation gives $Nx\in F$ for
every $x\in\mathbf Q$. Therefore
\[
 Nf(x)=f(Nx)=Nx.
\]
The nonzero rational scalar $N$ cancels, giving $f(x)=x$ for every $x$.
Lean checks this with the actual finite-index subgroup and
\texttt{nsmul\_index\_mem}. It does not assume that a named
``finitary'' predicate has already been proved trivial.
The identity satisfies all three conditions, so their union is exactly
the singleton $\{\mathrm{id}\}$. In particular the equivalence of
conditions 2 and 3 on this group is also a formal theorem.

\section{Any genuine metric fails to provide the claimed approximation}
First let $d$ be any metric on $\mathbf Q$, with no compatibility
assumption about its additive or field operations. Since $1\ne2$,
\[
 \varepsilon_0=d(1,2)>0.
\]
Every candidate $f$ in the union above equals the identity. At $x=1$,
\[
 d(f(1),D(1))=d(1,2)=\varepsilon_0.
\]
Consequently it is impossible to find a candidate with
$d(f(x),D(x))<\varepsilon_0$ for all $x$. The closed Lean theorems
negate the full universal approximation claim separately for literal
finite support and for the standard finite-index subgroup condition.

If the source instead places the metric directly on
$\operatorname{Aut}(\mathbf Q,+)$, let
$\varepsilon_1=d(\mathrm{id},D)>0$. Every candidate has distance
exactly $\varepsilon_1$ from $D$, again preventing arbitrary
approximation. This second interpretation is separately formalized.
No formula for an undefined distance is inserted into either proof.
More generally, for every $T_1$ topology on the actual automorphism
group, the candidate set is the closed singleton $\{\mathrm{id}\}$,
so $D$ is not in its closure. This separately proved topological
obstruction also covers genuine extended metrics that may take the
value infinity.

This argument does not treat a nonseparating pseudometric as a metric.
Nor does it identify finite rank relative to a chosen vector-space
basis with finite support of group automorphisms. That different
linear-algebra notion is not specified by the source and is not
claimed to be refuted here.

\section{Verification}
Lean 4.33.0 and Mathlib commit
\texttt{db584cd6d46c92f209a44c0f1c829460d327499d} are pinned.
The package includes a new root build, twenty-two transitive axiom audits
and full official \texttt{leanchecker --fresh Main} replay with
source-bound commands, exit codes, logs and hashes. No \texttt{sorry},
\texttt{native\_decide} or additional axiom is used. OpenAI Codex
assisted; a separate supervising Codex agent reviewed the source,
report and PDF. This is AI review, not human certification.
Organizer acceptance is pending.
\end{document}
